% ==========================================================
% == Section 1: Introduction                               ==
% ==========================================================
\section{Introduction}
\label{sec:intro}

% [Figure 1] 占位：用户后续将补充 motivation 图（左:MAGIC 攻击策略坍缩热力图;中:DACE 结构化策略覆盖;右:组合攻击涌现示例）。
% Placeholder figure — user will provide the actual PDF.
\begin{figure}[t]
  \centering
  \fbox{\rule[-3.3cm]{0pt}{6.6cm}\rule[-3.3cm]{0.95\linewidth}{0pt}}
  \caption{\textbf{Motivation and overview of \method.} \textbf{(a)} Existing co-evolutionary frameworks leave the attacker free to collapse to a narrow set of rewrite strategies, so the defender is trained on a biased slice of the threat space. \textbf{(b)} \method explicitly structures the attack space into a $12\times10$ risk-by-style grid, optimizes a normalized marginal coverage reward, and maintains a Bayesian adversarial replay pool that tracks per-sample threat posteriors, producing both \emph{broad} and \emph{persistent} co-evolution. \textbf{(c)} As training proceeds, \method's attacker exhibits \textit{combinatorial emergence}: single rewrites compose multiple risk categories and attack styles, a behavior not observed in prior co-evolution baselines.}
  \label{fig:motivation}
\end{figure}

% [段 1 · Hook] 大意:以攻击持续演化开篇,建立"威胁不断升级、静态对齐追不上"的紧迫感,呼应 MAGIC 开篇但不完全重复。
% 撰写逻辑:从模型部署的成功切入 → 攻击演化的速度（静态模板 → 自动化 → 组合多步 → agentic) → 由此点出静态对齐方法对"预先收集分布"的依赖带来的结构性局限 → 引出"必须用动态方法应对动态威胁"的转折。
% 中文对应内容:大语言模型已广泛部署在代码生成、复杂推理和 agent 系统中,然而安全事件频发。威胁形式正在从早期的人工越狱模板（DAN）,经由梯度优化（GCG)和 LLM 迭代改写（PAIR、TAP、AutoDAN）,演化到今天的多轮 agent 化攻击（X-Teaming、MAJIC）。这些攻击的共同特征是不断涌现新的策略组合。相应地,依赖预先收集分布进行一次性对齐的方法（RLHF、guardrail、静态 red-teaming 数据集)在面对分布漂移的攻击时会持续失守。这提示我们:对抗持续演化的威胁需要同样持续演化的防御,而不仅是一次更强的静态对齐。
Large language models (LLMs) have been rapidly deployed across high-stakes applications, from code and scientific reasoning~\citep{bai2025intern} to open-ended agent systems. Yet their safety alignment is repeatedly outpaced by an accelerating threat landscape: from hand-crafted jailbreak templates~\citep{shen2023doanythingnow} to gradient-based optimization~\citep{zou2023universal}, LLM-driven iterative rewriting~\citep{chao2023pair, mehrotra2024tap, liu2023autodan}, and, most recently, multi-turn agentic exploitations~\citep{rahman2025xteaming, qi2025majic, samvelyan2024rainbow}. Each wave introduces qualitatively new composition patterns that are absent from prior red-teaming corpora. Against such non-stationary adversaries, alignment pipelines that depend on a pre-collected distribution---be it RLHF preferences~\citep{ouyang2022training, bai2022constitutional, dai2023safe}, fixed guardrails, or static red-team datasets~\citep{jiang2024wildteaming, xie2024sorry}---are structurally ill-suited: they can harden the model only against the slice of attacks they have already seen. Confronting an attack distribution that \emph{keeps moving} thus demands a defense that \emph{keeps moving with it}.

% [段 2 · 协同演化范式 + 双病态] 大意:介绍协同演化（co-evolution）作为"用动态方法应对动态威胁"的最新范式,然后点出两个相互加强的核心病态。这是整篇论文的问题框架。
% 撰写逻辑:先肯定协同演化的最新进展（Self-RedTeam 的零和自博弈、MAGIC 的非对称序贯博弈、AdvEvo-MARL 的多智能体) → 转折指出"攻防都在进步"不等于"攻击覆盖足够广 + 防御记忆足够久" → 定义并命名两个病态:攻击策略坍缩（Attack Strategy Collapse）、防御对抗遗忘（Defense Adversarial Forgetting) → 解释二者相互加强的反馈回路（坍缩→偏科防御→巩固窄策略）。
% 中文对应内容:为应对动态攻击分布,研究者提出了一系列将攻击者和防御者置于同一 RL 博弈中的协同演化框架:Self-RedTeam 证明了零和自博弈下 Nash 均衡可带来安全保证;MAGIC 将交互建模为非对称序贯博弈并使用子博弈完美 Nash 均衡（SPNE）;AdvEvo-MARL、SSP 等进一步扩展了多智能体或自博弈变体。这些工作共同证明协同演化可以让对齐"跟上"攻击。但我们观察到,"都在进步"不等于"覆盖足够广、记忆足够久"。具体地,在真实 RL 训练中会涌现两个相互加强的病态:攻击策略坍缩——攻击者沿梯度过度强化少数高回报模式,攻击分布严重不均;防御对抗遗忘——当攻击者策略漂移到新区域时,防御者对旧攻击模式的拒答能力持续退化。更严重的是这两个病态互相加强:坍缩的攻击分布训出只擅长应付窄策略的防御者,而偏科的防御者又进一步把攻击者推回这些"容易得分"的策略区,形成恶性循环。
Recent work addresses this non-stationarity with \emph{co-evolutionary} training, in which an attacker LLM and a defender LLM are jointly optimized in an adversarial game~\citep{liu2025chasing, magic2025, pan2025advevo, ssp2026, triplayrl2026}. These frameworks show that under suitable equilibrium notions---Nash in zero-sum self-play, or subgame-perfect Nash in asymmetric sequential games---the defender can provably keep up with an evolving attacker. However, ``both sides keep improving'' is a weaker property than ``attacks cover the true threat space and defenses \emph{remain} robust to past threats.'' In practice we find that online RL-based co-evolution exhibits two compounding pathologies. First, \textbf{attack strategy collapse}: after the early exploration phase, the attacker overfits to a narrow set of high-reward rewrite patterns; the strategy distribution of its successful attacks becomes heavily skewed (far from uniform over the true policy space), so the defender is ultimately trained on a biased slice of the threat landscape~\citep{samvelyan2024rainbow, hong2024curiosity}. Second, \textbf{defense adversarial forgetting}: as the attacker drifts into new regions, the defender's competence on \emph{earlier} attack patterns silently decays~\citep{ssp2026}, mirroring catastrophic forgetting in continual learning. Worse, the two pathologies reinforce each other---a collapsed attack distribution trains a narrow defender, and a narrow defender pushes the attacker back toward the same easy regions---so standard co-evolution, left untreated, converges to a shallow equilibrium that underestimates the true safety gap.

% [段 3 · Gap 分析（关键差异化段落）] 大意:逐点分析现有工作"静态机制嫁接到动态系统"的两种具体失败模式——diversity 层面和 replay 层面——为 DACE 的两个核心组件做铺垫。
% 撰写逻辑:先拆 diversity:(i) 自动红队的 QD 流派针对静态目标,不适合动态协同;(ii) RL diversity 流派仅在文本/嵌入层面做多样性,与策略层无关;(iii) 特别点出 TriPlay-RL 是在协同演化中已注意到 diversity 的最接近先验工作,但其信号仅基于 Self-BLEU + cosine similarity,属于"换皮不换招"。再拆 replay:承认 SSP 确实实时更新样本风险（修正 opus/gpt 的误解）,但指出其更新依赖离散 safety score,遗漏三个建模维度（不确定性 / 防御者偶然性 / 防御者动态性）。最后引出 DACE 的两个核心设计方向。
% 中文对应内容:现有试图缓解上述病态的工作可以归纳为"将静态机制嫁接到动态系统"。在多样性一侧,自动红队领域的 quality-diversity 方法（Rainbow Teaming、QDRT、RainbowPlus、Ferret）针对静态目标模型构建 archive,不考虑防御者在联合训练中持续演化;RL 驱动的多样性奖励方法（CRT、DiveR-CT、GFlowNet）停留在文本或嵌入的语义层面,攻击者只需变换措辞即可骗过多样性惩罚。特别地,TriPlay-RL 是在协同演化建模中已明确关注多样性的最接近先验工作,但其多样性信号仅由 Self-BLEU 和句向量余弦距离构成,属于纯表层约束:攻击者仍可保持策略单调——"换皮不换招"。在回放/记忆一侧,SSP（Be Your Own Red Teamer)设置了经验回放池,并会在每次回放后依据 judge 给出的离散 safety score 实时更新样本风险;这已经是正确方向的第一步。但其离散更新机制仍然缺失三个关键维度:(i) 不表达估计的不确定性——采样少、估计不稳定的样本和被多次验证过的样本被一视同仁;(ii) 不考虑防御者响应的偶然性——单次拒答可能来自采样温度,离散机制可能据此过早淘汰真正危险的攻击;(iii) 不建模防御者能力的动态漂移——累计分数没有遗忘机制,旧的估计会持续污染对当前防御者的判断。要同时解决这两侧的不足,需要（a）在攻击策略空间而非文本空间显式做多样性,并且（b）用具备不确定性表达与时间衰减能力的连续值机制替代离散回放更新。
Existing attempts to mitigate these pathologies largely graft static mechanisms onto this inherently dynamic system, yielding two structurally incomplete directions. On the \emph{diversity} side, one family inherits the quality-diversity (QD) paradigm from automated red-teaming, constructing fixed behavior-descriptor archives~\citep{samvelyan2024rainbow, han2024ruby, pala2025ferret, dang2025rainbowplus, wang2025qdrt}; because the target is treated as static, these archives do not adapt to a defender that is itself under optimization. A second family attaches intrinsic-diversity rewards to RL attackers~\citep{hong2024curiosity, zhao2025diverct, lee2025gflownet, wang2025jailbreakr1, beutel2024openaidiverse}, but measures diversity at the \emph{textual} or \emph{embedding} level (e.g., Self-BLEU, cosine distance), so the attacker can satisfy the diversity signal merely by rephrasing while the underlying strategies stay collapsed. Most relevant to us, \textbf{TriPlay-RL}~\citep{triplayrl2026} is the closest prior work that explicitly injects diversity \emph{inside} a co-evolutionary attacker-defender-judge loop; yet its diversity term is a composition of Self-BLEU and sentence-embedding cosine similarity---i.e., still a purely semantic-surface constraint. Attackers can therefore pass the diversity filter by varying wording while reusing the same strategy: ``new skin, same trick.'' On the \emph{memory} side, SSP~\citep{ssp2026} correctly recognizes the forgetting problem and does maintain a replay pool whose sample-level harm scores are refreshed on each replay. Crucially, however, its update is driven by a \emph{discrete} safety score from the judge and therefore misses three modeling axes that are essential under a drifting defender: (i) it cannot express \emph{estimation uncertainty}---a sample tried once and a sample verified a hundred times contribute identically; (ii) it does not account for the \emph{stochasticity} of a single defender response---a one-shot refusal, possibly due to sampling temperature, can prematurely purge an attack that remains dangerous in expectation; and (iii) it has no \emph{forgetting} mechanism to track defender \emph{non-stationarity}---stale evidence keeps contaminating current threat estimates. These limitations motivate a co-evolutionary framework that places diversity at the \emph{strategy level} rather than the textual level, and that replaces discrete harm accumulation with a principled continuous posterior supporting uncertainty, stochasticity, and decay.

% [段 4 · DACE 提案 + "we observe" 钩子] 大意:以"In this paper, we propose DACE..."作为转折,简明介绍 DACE 的两大核心组件,并用一行"we observe"语句预告组合式攻击涌现这一强实验钩子（配合 Figure 1c）。
% 撰写逻辑:开头直接说我们提出 DACE → 一句话概括机制（结构化策略空间 + 归一化覆盖增益奖励 + Beta-Bernoulli 贝叶斯回放池) → 用一句话点出实验中涌现的"组合攻击"现象作为视觉钩子 → 指向 Figure 1。
% 中文对应内容:为此我们提出 DACE,一个以多样性驱动的协同演化框架,显式处理上述两大病态。一方面,DACE 在攻击者端引入来自 Llama-Guard-4 的 12 类风险和 Rainbow Teaming 的 10 种攻击风格共同构成的 12×10 策略空间,并通过归一化边际覆盖增益奖励驱动攻击探索;归一化消除了随档案池增长而衰减的奖励,并具有直接的 KL 散度收缩效率解读。另一方面,DACE 在防御者端构建统一的贝叶斯对抗回放池:每一条成功攻击维护一个 Beta-Bernoulli 后验,随防御者训练通过时间衰减自适应,并用 Thompson Sampling 在"利用已验证高危攻击"与"探索不确定性大的攻击"之间取得自然平衡,只在被充分验证后才被剪枝。我们在训练后期观察到一个 MAGIC 和 Self-RedTeam 未曾出现的现象——组合式攻击涌现:单条攻击 prompt 在文本中无缝融合多个风险类别与攻击风格;该行为是 DACE 覆盖压力和记忆压力共同作用下的自然产物,也是我们主要实验章节重点分析的对象之一。图 1 给出了整体示意。
In this paper, we propose \textbf{\method} (\textbf{D}iversity-Driven \textbf{A}dversarial \textbf{C}o-\textbf{E}volution), a co-evolutionary safety-alignment framework that jointly targets attack strategy collapse and defense adversarial forgetting. On the attacker side, \method exposes an explicit $12\times10$ attack strategy space---$12$ safety risk categories drawn from the Llama-Guard-4 taxonomy~\citep{llamaguard4_2026} and $10$ linguistic attack styles adapted from Rainbow Teaming~\citep{samvelyan2024rainbow}---and optimizes a \emph{normalized marginal coverage gain} reward whose denominator is the best single-sample entropy gain attainable at the current archive state. This normalization yields an $\mathcal{O}(1)$ training signal that does not vanish as the archive grows, and admits a closed-form interpretation as the fraction of one-step KL contraction (toward the uniform strategy prior) that a given rewrite achieves. On the defender side, \method maintains a single \emph{Bayesian adversarial replay pool} that serves dual roles: it tracks strategy frequencies for the coverage reward and stores every successful attack under a Beta--Bernoulli threat posterior. Time decay makes the posterior track defender drift, Thompson sampling selects replayed attacks that balance high-confidence threats (exploit) with uncertain yet promising ones (explore), and pruning is triggered only after repeated verification that a sample is genuinely neutralized. We observe a qualitatively new behavior in mid-to-late training that neither MAGIC nor Self-RedTeam reported: the \method attacker \emph{combinatorially composes} risk categories and attack styles within a single rewrite, a direct effect of the joint coverage and memory pressure (Figure~\ref{fig:motivation}).

% [Contributions] 大意:三条贡献 bullet,分别对应"问题框架 / 方法 / 实验",呼应 abstract 与 conclusion。
% 撰写逻辑:bullet 1 提出并命名两个病态,强调这是第一次把它们建模为联合问题;bullet 2 描述两个核心组件;bullet 3 实验证据。
% 中文对应内容:我们的贡献为:(1) 首次系统刻画协同演化中的双病态（攻击策略坍缩 + 防御对抗遗忘）,并论证二者需联合优化才能获得广度+持久度兼备的对齐;(2) 提出 DACE 框架,包括显式 12×10 攻击策略空间、归一化边际覆盖增益奖励与 Beta-Bernoulli 贝叶斯对抗回放池;(3) 在 9 个以上主流安全 benchmark 与 SOTA 自动红队工具上,DACE 在不牺牲通用能力的前提下取得更低的 ASR 与更广的策略覆盖,并首次观察到组合式攻击涌现。
\paragraph{Contributions.}
\begin{itemize}[leftmargin=1.2em, itemsep=2pt, topsep=2pt]
  \item \textbf{Problem framing.} We identify \emph{attack strategy collapse} and \emph{defense adversarial forgetting} as two compounding pathologies of RL-based co-evolutionary safety alignment, and argue that any framework seeking both \emph{broad} and \emph{persistent} robustness must treat coverage and memory as joint optimization objectives rather than independent heuristics.
  \item \textbf{Method.} We propose \method, which operationalizes this joint objective through an explicit $12\times10$ strategy space (a risk $\times$ style grid derived from Llama-Guard-4 and Rainbow Teaming), a normalized marginal coverage gain reward with a bounded non-vanishing signal and a closed-form KL-contraction interpretation, and a unified Bayesian adversarial replay pool with Beta--Bernoulli posteriors, time decay, and Thompson sampling.
  \item \textbf{Empirical.} On standard safety suites (WildGuardTest, HarmBench, DAN, XSTest, OR-Bench, StrongREJECT, WildJailbreak), strong-attacker OOD benchmarks (PAIR, TAP, AutoDAN-Turbo, MAJIC, Auto-RT, X-Teaming), and general-capability suites (IFEval, ARC-C, GPQA, MMLU, AlpacaEval~2), \method attains lower attack success rates than MAGIC, Self-RedTeam, and SSP without sacrificing utility, while its attacker achieves substantially broader strategy-space coverage and exhibits combinatorial attack emergence. Code and the strategy-indexed adversarial replay pool will be released.
\end{itemize}
